\section{Спецификация стенда: как его воспроизвести}
\label{sec:dishspec}

Ниже собрано всё, что нужно, чтобы собрать такой стенд заново: оборудование, петля, кодирование входа, чтение выхода, обратная связь и протокол опыта. Каждый параметр помечен источником. «Статья» --- значение взято из текста работы~\cite{Kagan2022}. «Выбор» --- в статье его нет, и для воспроизведения его придётся задать самому; мы предлагаем значение по умолчанию и говорим, как его проверить.

\subsection*{Оборудование и культура}

\begin{center}
\footnotesize
\setlength{\tabcolsep}{4pt}
\renewcommand{\arraystretch}{1.15}
\begin{tabular}{>{\raggedright}p{4.0cm}>{\raggedright}p{6.9cm}>{\raggedright\arraybackslash}p{1.6cm}}
\toprule
Параметр & Значение & Источник \\
\midrule
чип & матрица MaxOne: $26\,000$ платиновых электродов на площадке $8$~мм$^2$ & статья \\
запись & до $1024$ каналов одновременно, $20\,000$ отсчётов в секунду & статья \\
стимуляция & до $32$ электродов технически, в опыте --- $8$ независимых & статья \\
клетки & кора зародыша мыши; нейроны человека из стволовых клеток (два способа выращивания); клетки HEK293T (не нейроны) для контроля & статья \\
посев & около $10^6$ клеток на чип & статья \\
возраст культуры к опыту & мышь: с $\sim10$ суток; человек: $14$--$24$ суток или $\sim80$ суток в зависимости от способа выращивания & статья \\
\bottomrule
\end{tabular}
\end{center}

\subsection*{Петля и время}

Петля работает окнами по $10\ms$ ($200$ отсчётов): за окно считаются импульсы на электродах двигательных областей, игра обновляется, и выдаётся очередная стимуляция. Задержка от импульса в культуре до ответного стимула --- около $5\ms$. Во время игры сигнал каждого электрода проходит фильтр верхних частот Бесселя 2-го порядка с частотой среза $100$~Гц; модуль сигнала сглаживается фильтром нижних частот Бесселя 1-го порядка с частотой $1$~Гц, и порог импульса пропорционален этому сглаженному модулю. Порог в шесть стандартных отклонений фона статья указывает для отдельных измерений активности, а не для игры. Чтобы стимуляция не засчитывалась как ответ культуры, все сигналы гасятся, когда одновременно приходит больше $15$ крупных, выше $75$~мВ, импульсов. Всё это --- статья.

\subsection*{Кодирование мяча}

\begin{center}
\footnotesize
\setlength{\tabcolsep}{4pt}
\renewcommand{\arraystretch}{1.15}
\begin{tabular}{>{\raggedright}p{3.4cm}>{\raggedright}p{7.5cm}>{\raggedright\arraybackslash}p{1.6cm}}
\toprule
Параметр & Значение & Источник \\
\midrule
чувствительная область & $8$ электродов, расположенных так, чтобы соседние электроды означали соседние положения мяча & статья \\
код местом & номер электрода $k=1,\dots,8$ задаёт положение мяча относительно центра ракетки & статья \\
какая координата & вертикальное положение мяча; для воспроизведения --- относительно ракетки, $k=\lceil 8(y-p+\tfrac12)\rceil$ при $y-p\in[-\tfrac12,\tfrac12]$ & статья; формула --- выбор \\
код частотой & частота стимуляции от $4$ до $40$~имп/с несёт расстояние до ракетки & статья \\
направление шкалы & $4$~имп/с, когда мяч у противоположной стены, и линейно до $40$~имп/с, когда он у стены ракетки: $f=4+36\,(1-x/L)$, $x$ --- расстояние до стены ракетки, $L$ --- ширина корта & статья \\
форма импульса & короткий прямоугольный двухфазный: сначала положительная, потом отрицательная фаза & статья \\
амплитуда & $75$~мВ; выбрана так, чтобы не заставлять срабатывать заторможенные клетки & статья \\
длительность фаз & не указана; взять стандартную настройку системы стимуляции и не менять её в течение опыта & выбор \\
\bottomrule
\end{tabular}
\end{center}

\subsection*{Чтение ракетки}

В статье две двигательные области: импульсы в первой двигают ракетку вверх, во второй --- вниз; вклад областей взвешивается, чтобы учесть разную плотность клеток и разную активность~\cite{Kagan2022}. Точная формула не приведена. Для воспроизведения берём правило~\eqref{eq:dishreadout}, $\Delta p=c\,(n_\uparrow-n_\downarrow)$ за каждое окно $10\ms$, с весом, уравнивающим области по их средней активности в покое (выбор): $n_\uparrow$ и $n_\downarrow$ делятся на средний счёт своей области за окно, измеренный перед игрой. Коэффициент $c$ выбираем так, чтобы разница в один средний счёт сдвигала ракетку на $1\%$ высоты корта за окно (выбор); тогда постоянный перевес одной области проводит ракетку через весь корт за $1$~с.

Расположение областей на чипе в тексте статьи координатами не задано; есть только схема. Для воспроизведения достаточно трёх непересекающихся групп электродов: восемь чувствительных в центре и по группе записывающих электродов с двух сторон, одна --- «вверх», другая --- «вниз» (выбор).

\subsection*{Игра}

Размеры корта, длина ракетки и скорость мяча в статье не указаны; сказано только, что мяч летит с постоянной скоростью и отражается от стен. Для воспроизведения (выбор): корт $1\times1$, ракетка на правой стенке длиной $0.2$ (попадание при $|y-p|<0.1$, как в модели~\texttt{models/dishbrain\_model.py}), мяч пересекает корт за $1$~с. Промах без единого отбитого мяча в розыгрыше считается «эйсом»; розыгрыш длиннее трёх ударов подряд --- «длинным» (статья).

\subsection*{Обратная связь}

\begin{center}
\footnotesize
\setlength{\tabcolsep}{4pt}
\renewcommand{\arraystretch}{1.15}
\begin{tabular}{>{\raggedright}p{2.6cm}>{\raggedright}p{8.3cm}>{\raggedright\arraybackslash}p{1.6cm}}
\toprule
Событие & Что подаётся & Источник \\
\midrule
удар & предсказуемый стимул: все $8$ чувствительных электродов одновременно, $75$~мВ, $100$~имп/с, $100\ms$; обычная стимуляция мяча на это время прерывается & статья \\
промах & непредсказуемый стимул: $150$~мВ, $5$~имп/с, $4$~с, в случайные моменты на случайно выбранные из $8$ электродов; затем $4$~с без стимуляции, и мяч стартует в случайном направлении & статья \\
закон случайности & не указан; для воспроизведения --- моменты импульсов как пуассоновский поток со средней частотой $5$~имп/с & выбор \\
\bottomrule
\end{tabular}
\end{center}

\subsection*{Протокол и контроль}

Один сеанс длится $20$~мин; сравнивают первые $5$ и последние $15$~мин (статья). Три меры результата: средняя длина розыгрыша, доля эйсов и доля длинных розыгрышей. Условия опыта (статья):
\begin{itemize}
\item \emph{стимул} --- полная петля, как описано выше;
\item \emph{тишина} --- вместо стимула удара или промаха столько же времени без всякой стимуляции, затем мяч стартует в случайном направлении;
\item \emph{без обратной связи} --- мяч после промаха просто отражается и летит дальше, стимуляция положения мяча продолжается;
\item \emph{покой} --- игра идёт и ракетка движется по активности, но стимуляции нет совсем;
\item \emph{контроли} --- чип только со средой; клетки HEK293T; «культура в компьютере», симуляция вместо клеток.
\end{itemize}
Объём выборки в основной серии: $9$ культур мыши ($101$ сеанс), $11$ культур человека ($138$), $6$ чипов со средой ($80$), $20$ культур в покое ($42$), $3$ симуляции ($38$), всего $399$ сеансов. В серии про обратную связь --- $486$ сеансов за $3$ дня по $3$ сеанса в день, условия «стимул», «тишина» и «без обратной связи» ($n=27$, $17$ и $15$). Рост средней длины розыгрыша от первых $5$ минут к последним $15$ получен только у живых культур и только в условии «стимул»~\cite{Kagan2022}.

\subsection*{Цикл стенда по шагам}

Каждые $10\ms$ компьютер стенда делает следующее.
\begin{enumerate}
\item Считывает число импульсов $n_\uparrow$, $n_\downarrow$ в двух двигательных областях за прошедшее окно и нормирует их на средний счёт области в покое.
\item Сдвигает ракетку на $\Delta p=c\,(n_\uparrow-n_\downarrow)$ и ограничивает её краями корта.
\item Сдвигает мяч на один шаг; отражает его от верхней, нижней и левой стенок.
\item Если мяч дошёл до правой стенки: при $|y-p|<0.1$ --- удар, счёт розыгрыша растёт на один, подаётся стимул удара ($100\ms$); иначе --- промах, розыгрыш записывается, подаётся стимул промаха ($4$~с), затем $4$~с паузы, и мяч стартует заново.
\item Иначе выбирает электрод $k$ по вертикальному положению мяча и частоту $f$ по расстоянию до стены ракетки и подаёт на электрод $k$ двухфазные импульсы $75$~мВ с частотой $f$ до следующего окна.
\end{enumerate}
Этого достаточно, чтобы собрать стенд, совместимый с опубликованным, и проверить главный вопрос главы: учит ли культуру предсказуемость или простая разница между ответом на удар и на промах. Для этого к трём условиям статьи стоит добавить четвёртое, которого в ней нет: после промаха --- \emph{предсказуемый} стимул той же длительности и энергии, что и непредсказуемый. Если культура учится и так, дело в разнице, а не в предсказуемости.
